\chapter*{General Introduction}
\addcontentsline{toc}{chapter}{General Introduction}
\thispagestyle{MyStyle}

The deployment of fiber optic infrastructure across France, driven by the \textit{Plan France Tr\`es Haut D\'ebit} and ARCEP's regulatory framework, has created a large-scale FTTX/GPON deployment landscape. Orange France manages millions of premise connections through a hierarchical OLT/NRO/SRO/PAR/subscriber-terminal chain, whose administrative workflows form the daily reality of the \textit{Guichet OI: PAR France} unit.

The unit's fifteen consultants process roughly 112 dossiers per day across nine disconnected tools (Agilis, SharePoint, Microsoft Planner, and others) sharing no common data layer. An estimated 55 percent of consultant time goes to deterministic, rule-based, yet entirely manual tasks: POI validation takes fifteen to twenty minutes per email; CAFF returns received after hours go unrouted, breaching the one-working-day contractual KPI; daily distribution consumes thirty to forty-five minutes with no cross-day fairness mechanism; and LOC PAR completeness verification needs fifteen to thirty minutes of manual reading and field extraction per dossier.

This Final Year Project, conducted within Amaris Consulting and embedded in the Guichet OI: PAR France unit, was born from twenty-five weeks of dual fieldwork as both operational consultants and system architects. The result is the \textbf{FiberGate Ecosystem}, an integrated three-pillar architecture combining robotic process automation, multimodal document intelligence, and a governed enterprise web platform.

The \textbf{SENTINEL Automation Fleet} comprises three Python agents built on the \texttt{exchangelib} EWS/NTLM library, orchestrated through a FastAPI asynchronous execution layer, and integrated with Microsoft 365 via authenticated Power Automate webhooks. Agent~1 automates POI validation through six-field regular-expression extraction dispatched to a seven-step SharePoint SUIVI/Planner flow; Agent~2 automates CAFF return routing through single-pattern OEIE identification and an atomic Power Automate flow guaranteeing immediate assignment; Agent~3 automates daily dossier distribution through a continuity-first round-robin pandas pipeline producing a three-sheet \texttt{openpyxl} workbook and per-task webhooks.

The \textbf{FiberGate AI Service} addresses the LOC PAR completeness verification workflow with a LayoutLMv3 multimodal transformer (Microsoft Research, 2022) fine-tuned through a two-head strategy combining six-class sequence classification and BIO token classification for twelve-field structured extraction. A confidence-gated policy routes uncertain extractions to human review, preserving audit integrity in this regulated environment.

The \textbf{Unified Platform} integrates both pillars in an Angular~17/Spring Boot~3.2.4/PostgreSQL stack with JWT-secured RBAC and a FastAPI proxy for the agents' background jobs, deployed on Azure Container Apps via GitHub Actions/Azure DevOps CI/CD with blue-green releases.

All figures and tables are the authors' own work unless otherwise attributed. The report unfolds across six chapters: context and Sure Step methodology (\textbf{Ch.~1}); existing-system analysis and state of the art (\textbf{Ch.~2}); UML-based system design (\textbf{Ch.~3}); SENTINEL Fleet implementation (\textbf{Ch.~4}); the FiberGate AI Service and LayoutLMv3 fine-tuning (\textbf{Ch.~5}); and the Unified Platform and its impact assessment (\textbf{Ch.~6}).
